\chapter{Réseaux pour le trading}\label{ch:nw:networking-for-trading}

L'opérateur du flux consolidé des options américaines indique à ses abonnés quelle capacité réseau acheter, et il ne donne pas une moyenne.
Pour juillet 2026, il projette 4.403 gigabits sur les 100 millisecondes les plus chargées et 0.501 gigabit sur les 10
millisecondes les plus chargées, par copie du flux~: \qty{44}{\giga\bit\per\second} sur la fenêtre la plus longue, \qty{50}{\giga\bit\per\second}
sur la plus courte, soit le double pour un abonné qui prend les deux copies redondantes. Ses propres statistiques disent la même chose vu de l'autre côté~: la
milliseconde la plus chargée de juillet 2026 a transporté des messages à un débit cinq fois supérieur à celui de la seconde la plus chargée. Un réseau
dimensionné pour la moyenne d'une journée de trading perd des paquets dans ses millisecondes les plus chargées, et ce sont ces millisecondes-là que
une société de trading est payée pour voir.

Ce chapitre ouvre la première partie du livre, consacrée au réseau entre une plateforme de négociation et les serveurs de la société. Il chiffre le coût d'un octet sur le fil,
montre comment un paquet atteint de nombreux abonnés, explique pourquoi les données de marché circulent en datagrammes et les ordres en flux, et se termine
là où naissent la plupart des pertes~: sur un port de commutateur qui reçoit davantage qu'il ne peut émettre.

\section{La trame Ethernet et ce que coûte un octet sur le fil}

\begin{definition}[Carte réseau, trame Ethernet]\label{def:nw:networking-for-trading:frame}
Une \emph{carte réseau}\index{carte réseau} (NIC) est le dispositif qui relie un ordinateur à une liaison réseau~: 
elle transforme les bits présents sur le fil en données dans la mémoire de l'hôte, et inversement. Une \emph{trame Ethernet}\index{trame Ethernet} est
l'unité qu'une liaison Ethernet transporte~: une adresse de destination et une adresse source (6 octets chacune), un type (2 octets), une charge utile de 46 à 1\,500
octets, une séquence de contrôle de trame (4 octets), soit de 64 à 1\,518 octets au total sans étiquette VLAN. Sur le fil, chaque trame est
précédée de 8 octets de préambule et de délimiteur de début et suivie d'au moins 12 octets d'intervalle inter-trames inactif.
\end{definition}

\begin{definition}[Débit de ligne, délai de sérialisation, unité de transmission maximale]\label{def:nw:networking-for-trading:ser}
Le \emph{débit de ligne}\index{débit de ligne} (line rate) $R_{\mathrm L}$ d'une liaison est le débit auquel elle cadence les bits sur le support (1, 10, 25
ou \qty{100}{\giga\bit\per\second}). Le \emph{délai de sérialisation}\index{délai de sérialisation} d'une trame est le temps que met la liaison
pour la placer sur le fil. L'\emph{unité de transmission maximale}\index{unité de transmission maximale} (MTU) est la charge utile la plus grande qu'une
liaison transporte dans une seule trame, soit 1\,500 octets pour l'Ethernet standard.
\end{definition}

\begin{proposition}[Coût d'une trame]\label{prop:nw:networking-for-trading:ser}
Une trame de $\ell_{\mathrm f}$ octets occupe $\ell_{\mathrm f} + 20$ octets de temps de ligne, si bien que son \omterm{def:nw:networking-for-trading:ser}{délai de sérialisation} vaut
\[
  t_{\mathrm{ser}} = \frac{8\,(\ell_{\mathrm f} + 20)}{R_{\mathrm L}},
\]
et une liaison transporte au plus $R_{\mathrm L} / (8(\ell_{\mathrm f}+20))$ trames par seconde. Tout dispositif qui doit recevoir une trame entière
avant de la réémettre paie $t_{\mathrm{ser}}$ une fois de plus.
\end{proposition}

\begin{proof}
Le préambule, le délimiteur de début et l'intervalle inter-trames représentent 20 octets de temps de ligne qui appartiennent à chaque trame et à aucune autre~;
une liaison émet un bit toutes les $1/R_{\mathrm L}$. Un dispositif de type store-and-forward (chapitre~2) ne commence à émettre qu'une fois le dernier bit
arrivé, soit un retard d'une sérialisation au débit entrant.
\end{proof}

\begin{omfigure}
\begin{tikzpicture}[font=\scriptsize,f/.style={draw,minimum height=0.95cm,minimum width=1.32cm,text width=1.2cm,
  align=center,inner sep=1pt}]
\foreach \i/\n/\b/\c in {0/{préam\-bule\\+ SFD}/8/gray!15, 1/en-tête\\Ethernet/14/omDef!15, 2/en-tête\\IPv4/20/omDef!15,
  3/en-tête\\UDP/8/omDef!15, 4/{Mold-\\UDP64}/20/omMeth!18, 5/lon\-gueur/2/omMeth!18, 6/ajout\\d'ordre/36/omThm!25,
  7/FCS/4/omDef!15, 8/{inter\-valle}/12/gray!15}
  {\node[f,fill=\c,anchor=south west] at (1.35*\i,0) {\n}; \node[font=\scriptsize] at (1.35*\i+0.66,-0.25) {\b\ octets};}
\draw[decorate,decoration={brace,amplitude=4pt,mirror}] (1.35,-0.5) -- (10.8,-0.5)
  node[midway,below=4pt,font=\scriptsize] {Trame Ethernet~: 104 octets (64 au minimum, 1\,518 au maximum)};
\draw[decorate,decoration={brace,amplitude=4pt}] (0,1.05) -- (12.15,1.05)
  node[midway,above=4pt,font=\scriptsize] {temps de ligne~: 124 octets $=$ \qty{99.2}{\nano\second} à \qty{10}{\giga\bit\per\second},
  \qty{39.7}{\nano\second} à \qty{25}{\giga\bit\per\second}};
\end{tikzpicture}

\omcaption{Un message d'ajout d'ordre sur le fil, au format de flux du simulateur (un paquet MoldUDP64 dans un datagramme UDP dans une
\omterm{def:nw:networking-for-trading:frame}{trame Ethernet}). Les 36 octets que lit la stratégie représentent 29~\% du temps de ligne qu'occupe la trame. Les cadres ne sont pas à l'échelle~; les nombres d'octets
proviennent de la \cref{prop:nw:networking-for-trading:ser} et des en-têtes des formats.}\label{fig:nw:networking-for-trading:frame}
\end{omfigure}

\begin{example}[Un ajout d'ordre]\label{ex:nw:networking-for-trading:mold}
Le simulateur de bourse du livre~10 publie son flux sous forme de paquets MoldUDP64~: un en-tête de 20 octets (session, numéro de
séquence, nombre de messages), puis chaque message précédé d'une longueur de 2 octets. Un message d'ajout d'ordre fait 36 octets, si bien qu'un paquet qui en transporte un seul
fait 58 octets de charge utile UDP~; avec 8 octets d'en-tête UDP, 20 d'en-tête IPv4 et 18 d'en-tête Ethernet et de séquence de contrôle, la trame
fait 104 octets et occupe 124 octets de temps de ligne (\cref{fig:nw:networking-for-trading:frame})~: \qty{99.2}{\nano\second}
à \qty{10}{\giga\bit\per\second}, \qty{39.7}{\nano\second} à 25. Une plateforme qui regroupe huit de ces messages dans un seul paquet
ne paie les 66 octets de surcharge fixe qu'une fois pour huit messages, au lieu de huit fois.
\end{example}

\begin{table}[ht]
\centering\small
\begin{tabular}{rrrrr}
\toprule
Trame (octets) & \qty{1}{\giga\bit\per\second} & \qty{10}{\giga\bit\per\second} & \qty{25}{\giga\bit\per\second} &
\qty{100}{\giga\bit\per\second} \\
\midrule
64 & 672.0 & 67.2 & 26.9 & 6.7 \\
128 & 1\,184.0 & 118.4 & 47.4 & 11.8 \\
256 & 2\,208.0 & 220.8 & 88.3 & 22.1 \\
512 & 4\,256.0 & 425.6 & 170.2 & 42.6 \\
1\,024 & 8\,352.0 & 835.2 & 334.1 & 83.5 \\
1\,518 & 12\,304.0 & 1\,230.4 & 492.2 & 123.0 \\
\bottomrule
\end{tabular}
\caption{\omterm{def:nw:networking-for-trading:ser}{Délai de sérialisation} en nanosecondes selon la taille de la trame et le \omterm{def:nw:networking-for-trading:ser}{débit de ligne} (\cref{prop:nw:networking-for-trading:ser}, préambule
et intervalle inter-trames inclus). Calculé par \texttt{nw\_wire.frame\_table}.}\label{tab:nw:networking-for-trading:ser}
\end{table}

Le \Cref{tab:nw:networking-for-trading:ser} est le premier tableau de latence de ce livre et celui que l'on oublie le plus souvent. Une
trame de taille maximale à \qty{10}{\giga\bit\per\second} met \qty{1.23}{\micro\second} à se sérialiser, soit la moitié des
\qty{2.5}{\micro\second} qu'un ingénieur d'un grand teneur de marché (market maker) a présentés lors d'une conférence publique comme un très bon temps de bout en bout pour un
système de trading logiciel complet (livre~13, chapitre~1). Passer de 10 à \qty{25}{\giga\bit\per\second} fait gagner 60~\% sur
chaque sérialisation, ce qui justifie d'acheter le port le plus rapide même lorsque le port plus lent suffirait à écouler le trafic.
La sérialisation est aussi la raison pour laquelle les petites trames comptent~: un ordre dépasse rarement une centaine d'octets, et un commutateur qui le réémet
seulement après l'avoir reçu en entier en paie chaque octet.

\begin{remark}[Bits en vol]\label{rem:nw:networking-for-trading:flight}
La lumière traverse une fibre monomode standard à raison d'un mètre en \qty{4.88}{\nano\second} environ (indice de groupe 1.462, chapitre~10).
Une trame de 1\,518 octets à \qty{10}{\giga\bit\per\second} mesure donc 252~mètres dans la fibre, davantage que tout câble présent dans une
salle de centre de données~: son premier bit atteint l'extrémité lointaine bien avant que son dernier bit n'ait quitté la source. La distance et la sérialisation sont deux
termes distincts d'un budget de latence, et à l'échelle d'un bâtiment c'est généralement la sérialisation qui l'emporte.
\end{remark}

\omcode{code/firm/netsim/firm_netsim.py}{36}{50}{Calcul de la taille des trames dans \texttt{firm.netsim}~: la plus petite trame est complétée à 64
octets, et chaque trame paie 20 octets de préambule et d'intervalle.}\label{lst:nw:networking-for-trading:frame}

\section{Multicast IP~: groupes, appartenance et snooping}

Une plateforme de négociation envoie les mêmes données de marché à des centaines d'abonnés. Envoyer une copie à chacun sur sa propre connexion ferait
attendre le dernier abonné que tous les autres soient servis (le chapitre~12 étudie une plateforme qui a fait exactement cela)~; envoyer une seule copie que le réseau réplique
les rend tous égaux, à la réplication près.

\begin{definition}[Groupe multicast, IGMP]\label{def:nw:networking-for-trading:group}
Un \emph{groupe multicast}\index{groupe multicast} est une adresse IPv4 comprise entre 224.0.0.0 et 239.255.255.255 qui désigne un ensemble
de récepteurs plutôt qu'un seul hôte~: un datagramme qui lui est envoyé est remis à chacun de ses membres. L'\emph{Internet Group Management
Protocol}\index{Internet Group Management Protocol} (IGMP) est le protocole par lequel un hôte indique à ses routeurs voisins
les groupes qu'il souhaite recevoir et, dans la version~3, les sources dont il accepte les flux.
\end{definition}

Le flux d'une plateforme de négociation est en général réparti sur de nombreux groupes, un par canal ou par partition d'instruments,
afin qu'un abonné ne reçoive que les groupes dont il a besoin (le chapitre~22 dimensionne ainsi les flux d'options). Un groupe doit aussi être livré sur le segment Ethernet,
et les adresses Ethernet comptent 48 bits~: une adresse de groupe est associée à l'adresse Ethernet \texttt{01:00:5e} suivie des 23 bits de poids faible du groupe,
si bien que les 5 bits de poids fort parmi les 28 bits qui distinguent les groupes sont perdus et que 32 groupes partagent chaque adresse Ethernet.
Une carte qui filtre sur les adresses Ethernet peut donc laisser passer des groupes auxquels l'hôte ne s'est jamais joint,
et le noyau les rejette après en avoir payé le coût.

\begin{definition}[IGMP snooping]\label{def:nw:networking-for-trading:snooping}
L'\emph{IGMP snooping}\index{IGMP snooping} est la pratique d'un commutateur qui lit les rapports d'appartenance IGMP qui le traversent et
n'achemine le trafic de chaque groupe que vers les ports sur lesquels un membre s'est signalé, au lieu d'inonder
chaque port de chaque trame multicast.
\end{definition}

\begin{omfigure}
\begin{tikzpicture}[font=\footnotesize,b/.style={draw,rounded corners=2pt,minimum height=0.65cm,align=center,inner sep=3pt},
  >=Stealth]
\node[b,fill=gray!10] (v) at (0,0) {flux de la plateforme\\ groupes A, B};
\node[b,fill=omDef!15,minimum width=1.9cm,minimum height=3.6cm] (s) at (3.6,0) {commutateur\\ (IGMP\\ snooping)};
\foreach \y/\t/\g in {1.35/{serveur 1~: abonné à A}/{A}, 0.45/{serveur 2~: abonné à A et B}/{A et B}, -0.45/{serveur 3~: abonné à B}/{B},
  -1.35/{serveur 4~: abonné à rien}/{rien}}
  {\node[b,minimum width=4.1cm,anchor=west] (h) at (7.4,\y) {\t};
   \draw[->,thick] (s.east|-h.west) -- node[above,font=\scriptsize]{envoie \g} (h.west);}
\draw[->,thick] (v) -- node[above,font=\scriptsize]{A et B} (s.west|-v.east);
\end{tikzpicture}

\omcaption{\omterm{def:nw:networking-for-trading:snooping}{IGMP snooping}. Le commutateur apprend, à partir des rapports d'appartenance des serveurs, quels ports veulent quel groupe, et réplique
chaque groupe uniquement vers eux~; le serveur~4 ne reçoit rien. Sans snooping, les quatre ports transporteraient les deux groupes, et les serveurs qui ne sont abonnés à rien
consacreraient tout de même du temps de leurs cartes et de leur noyau à les rejeter.}\label{fig:nw:networking-for-trading:snoop}
\end{omfigure}

\begin{example}[Deux groupes, une adresse Ethernet]\label{ex:nw:networking-for-trading:mac} \sloppy
Le groupe 239.255.1.2 est associé à \texttt{01:00:5e:7f:01:02}~: les 23 bits de poids faible de 255.1.2 valent 127.1.2 une fois le bit de poids fort de 255 retiré. Le groupe 224.127.1.2 est associé à la même adresse.
Un serveur qui ne s'est joint qu'au premier reçoit aussi, depuis tout commutateur ou toute carte qui filtre sur la seule adresse Ethernet,
le trafic du second. Les plateformes attribuent leurs groupes en tenant compte de cela, et
une société de trading qui exploite son propre multicast interne (le ticker plant du chapitre~27) devrait faire de même.
\end{example}

\section{Datagrammes contre flux}

Deux protocoles de transport acheminent presque tout ce qui circule dans un réseau de trading. UDP envoie des datagrammes~: chacun est remis entier
ou pas du tout, sans ordre ni retransmission, et un datagramme envoyé à un groupe atteint chacun de ses membres. TCP achemine un flux d'octets
entre deux extrémités~: il numérote les octets, les acquitte, retransmet ceux qui sont perdus
et livre le tout dans l'ordre.

Les données de marché circulent en multidiffusion UDP parce que seul le multicast envoie une copie à tous, et parce qu'un flux ne doit pas ralentir
pour son abonné le plus lent~: le contrôle de flux de TCP ferait exactement cela. Le prix à payer est que la perte devient le problème de l'application. La plateforme
envoie chaque paquet sur deux lignes indépendantes (livre~10, chapitre~26), les numérote et exécute des services de retransmission
et d'instantanés~; le feed handler de la société arbitre et récupère (livre~13, chapitres~16 et~18). Les ordres voyagent sur TCP parce qu'un ordre doit arriver
exactement une fois, dans l'ordre par rapport aux autres ordres de la société, et parce que la plateforme doit savoir quelle connexion l'a envoyé~: la couche session
du protocole de saisie des ordres repose sur ces garanties.

\begin{definition}[Blocage en tête de file]\label{def:nw:networking-for-trading:hol}
Le \emph{blocage en tête de file}\index{blocage en tête de file} est le retard imposé à chaque élément d'une file ordonnée par un élément en tête
qui ne peut pas avancer. Dans TCP, un segment perdu retient chaque octet ultérieur du flux, déjà reçu ou non, jusqu'à
l'arrivée de sa retransmission.
\end{definition}

Le \omterm{def:nw:networking-for-trading:hol}{blocage en tête de file} est le prix de l'ordre imposé par TCP. Une annulation envoyée une microseconde après un nouvel ordre perdu attend la retransmission
de ce nouvel ordre, qui arrive au moins après un aller-retour et, en général, après une temporisation~: sur une liaison de colocation sans perte, la perte est rare,
et lorsqu'elle survient, toute la session se bloque. L'ordre est aussi une qualité~: la plateforme ne doit jamais voir l'annulation avant
l'ordre qu'elle annule.

\begin{definition}[Algorithme de Nagle]\label{def:nw:networking-for-trading:nagle}
L'\emph{algorithme de Nagle}\index{algorithme de Nagle} est la règle de TCP selon laquelle un émetteur retient de nouvelles données, au lieu de les envoyer
dans un petit segment, tant que les données qu'il a déjà envoyées restent non acquittées. Il économise de la bande passante lorsqu'un programme écrit de nombreux petits morceaux,
et il a été conçu pour les terminaux interactifs.
\end{definition}

\begin{proposition}[Nagle rencontre l'acquittement différé]\label{prop:nw:networking-for-trading:nagle}
Supposons qu'un émetteur écrive un petit ordre puis, avant son acquittement, un second petit ordre, sur une connexion où l'\omterm{def:nw:networking-for-trading:nagle}{algorithme de Nagle} est activé,
vers un récepteur qui retarde ses acquittements (TCP autorise jusqu'à une demi-seconde). Le second ordre ne part que lorsque le premier est acquitté~:
après la temporisation d'acquittement différé du récepteur, ou après un aller-retour si le récepteur a lui-même des données à envoyer.
Désactiver l'\omterm{def:nw:networking-for-trading:nagle}{algorithme de Nagle} (l'option de socket \texttt{TCP\_NODELAY}) supprime l'attente.
\end{proposition}

\begin{proof}
Par la règle, la seconde écriture est retenue tant que le premier segment n'est pas acquitté. Le récepteur envoie son acquittement
lorsque sa temporisation expire, lorsqu'il dispose d'un second segment complet à acquitter (ce n'est pas le cas~: un seul petit segment est arrivé), ou avec ses propres données.
\end{proof}

\begin{method}[Un socket de saisie des ordres]\label{met:nw:networking-for-trading:socket}
\begin{enumerate}
\item Activer \texttt{TCP\_NODELAY} sur chaque connexion de saisie des ordres et de drop copy.
\item Écrire chaque message en un seul appel, de sorte qu'un message soit un seul segment et que rien n'attende une écriture ultérieure.
\item Lorsque le protocole de session de la plateforme acquitte au niveau applicatif, acquitter aussi rapidement au niveau TCP
  (l'option \texttt{TCP\_QUICKACK} de Linux n'est pas permanente et doit être réactivée après chaque lecture).
\item Garder la connexion ouverte (les heartbeats de la couche session) et préchauffée~: une nouvelle connexion TCP coûte un aller-retour avant
  son premier octet, et le chapitre~20 ajoute une poignée de main TLS (handshake) par-dessus pour les plateformes de crypto-actifs.
\end{enumerate}
\end{method}

\section{Congestion~: files de sortie, incast et micro-rafales}

Un commutateur achemine chaque trame vers un port de sortie. Si les trames destinées à un même port arrivent plus vite que le \omterm{def:nw:networking-for-trading:ser}{débit de ligne} de ce port,
elles attendent.

\begin{definition}[File de sortie, taux de sursouscription]\label{def:nw:networking-for-trading:egress}
Une \emph{file de sortie}\index{file de sortie} est la mémoire tampon dans laquelle un commutateur retient les trames en attente d'émission sur un port de sortie~;
lorsqu'elle est pleine, les trames qui arrivent sont rejetées (tail drop). Le \emph{taux de sursouscription}\index{taux de sursouscription} d'un port est la somme des \omterm{def:nw:networking-for-trading:ser}{débits de ligne} pouvant émettre vers lui,
divisée par son propre \omterm{def:nw:networking-for-trading:ser}{débit de ligne}.
\end{definition}

\begin{definition}[Incast, micro-rafale]\label{def:nw:networking-for-trading:burst}
L'\emph{incast}\index{incast} est la congestion qui survient lorsque de nombreux émetteurs transmettent simultanément vers un même récepteur. Un
\emph{micro-rafale}\index{micro-rafale} (microburst) est une période de quelques microsecondes à quelques millisecondes pendant laquelle le trafic offert à un port
dépasse son \omterm{def:nw:networking-for-trading:ser}{débit de ligne}, alors que sa moyenne sur n'importe quel intervalle qu'un outil de supervision rapporte (une seconde, une minute) reste bien en deçà
de celui-ci.
\end{definition}

Un flux qui occupe en moyenne la moitié d'un port peut le saturer pendant une milliseconde, et un réseau de trading est conçu pour que cela
se produise~: le flux de chaque plateforme, et chaque partition de chaque flux, réagit aux mêmes événements de marché au même moment.

\begin{proposition}[La borne fluide]\label{prop:nw:networking-for-trading:fluid}
Soient $n$ entrées qui émettent aux débits $r_1, \dots, r_n$ vers un port de \omterm{def:nw:networking-for-trading:ser}{débit de ligne} $R_{\mathrm L}$ pendant une rafale de durée $D$, à partir
d'une file vide. À la fin de la rafale, la file contient
\[
  B = \Bigl(\textstyle\sum_i r_i - R_{\mathrm L}\Bigr)^{+} D / 8 \text{ octets},
\]
et le dernier octet de la rafale quitte le port $8B/R_{\mathrm L}$ après son arrivée. Le rejet en queue ne perd rien si et seulement si la mémoire
disponible pour le port est au moins égale à $B$~; sinon il perd $B$ moins la mémoire, par trames entières.
\end{proposition}

\begin{proof}
Pendant la rafale, la file croît au rythme de la différence des débits lorsque celle-ci est positive, et reste vide dans le cas contraire~; le port
la vide ensuite à $R_{\mathrm L}$. La version discrète, trame par trame, est la boucle de la \cref{lst:nw:networking-for-trading:egress},
qui se réduit à la formule fluide lorsque les trames sont petites devant $B$.
\end{proof}

\begin{omfigure}
\begin{tikzpicture}
\begin{axis}[width=11.5cm,height=5.4cm,xmode=log,xmin=0.7,xmax=1400,ymin=0,ymax=380,
  xtick={1,10,100,1000},xticklabels={1~ms,10~ms,100~ms,1~s},
  xlabel={fenêtre de moyennage (échelle logarithmique)},ylabel={débit de crête (millions de messages / s)},
  tick label style={font=\small},label style={font=\small},grid=major,grid style={gray!25},
  legend columns=3,legend style={font=\footnotesize,at={(0.5,-0.32)},anchor=north,draw=none,fill=none,
  /tikz/every even column/.append style={column sep=8pt}},table/col sep=comma]
\addplot[omExo,very thick,mark=*,mark size=1.6pt] table[x=window_ms,y=2024_01]{figdata/networks/01-networking-for-trading/opra.csv};
\addplot[omMeth,very thick,mark=square*,mark size=1.6pt] table[x=window_ms,y=2025_03]{figdata/networks/01-networking-for-trading/opra.csv};
\addplot[omThm,very thick,mark=triangle*,mark size=2pt] table[x=window_ms,y=2026_07]{figdata/networks/01-networking-for-trading/opra.csv};
\legend{janv. 2024,mars 2025,juil. 2026}
\end{axis}
\end{tikzpicture}

\omcaption{La fenêtre la plus chargée de chaque longueur pour le flux consolidé d'options américaines, exprimée en débit~: le nombre de pointes de messages en une
seconde, 100~ms, 10~ms et 1~ms, divisé par la fenêtre. Plus la fenêtre est courte, plus le débit est élevé~; en juillet 2026, la milliseconde la plus chargée
a atteint 5.5 fois la seconde la plus chargée. Données~: OPRA, \emph{Key Operating Metrics} (août 2026), transcrites dans
\texttt{data/networks/opra\_metrics.csv}.}\label{fig:nw:networking-for-trading:opra}
\end{omfigure}

La \cref{fig:nw:networking-for-trading:opra} est la version publiée de la \omterm{def:nw:networking-for-trading:burst}{micro-rafale}~: le même flux, au même mois, représente 63.9 millions de messages par seconde
ou 350 millions selon la fenêtre dans laquelle on le lit. L'avis de capacité de l'opérateur, que les abonnés utilisent pour dimensionner leurs liens, est
rédigé en pointes de 10 millisecondes pour cette raison
(\cref{dat:nw:networking-for-trading:opra})~; un port de commutateur ne moyenne pas sur 10 millisecondes, il met en file trame par trame.

\begin{dated}{2026-09}{Combien de réseau le flux d'options demande}\label{dat:nw:networking-for-trading:opra}
La projection de capacité de la SIAC pour OPRA en vigueur en juillet 2026 donne, pour une copie du flux, 13.575 millions de messages, 4.403
gigabits et 1.564 million de paquets dans les 100 millisecondes les plus chargées, et 1.562 million de messages, 0.501 gigabit et 169\,000
paquets dans les 10 millisecondes les plus chargées. Cela représente \qty{44.0}{\giga\bit\per\second} et \qty{50.1}{\giga\bit\per\second}~; environ 352
octets et 8.7 messages par paquet sur la fenêtre de 100 millisecondes. Les abonnés qui prennent les deux copies redondantes doublent la
bande passante, et l'avis demande 10~\% de plus pour les retransmissions. L'avis précise que l'intervalle de 10 millisecondes «~reflète le taux d'utilisation du système
pendant les rafales de trafic~».
\end{dated}

Pour voir le mécanisme plutôt que son empreinte, le tutoriel du chapitre simule huit flux vers un seul port de
\qty{10}{\giga\bit\per\second} avec \texttt{firm.netsim}. Chaque flux a une moyenne de 400\,000 paquets par seconde, en grappes (un événement de marché produit une suite de messages),
d'un à trente messages d'ajout d'ordre par paquet~; les huit ensemble offrent environ
\qty{5.2}{\giga\bit\per\second}, soit la moitié du port. Dans une première version, les grappes de chaque flux sont indépendantes. Dans l'autre, 30~\% des
paquets de chaque flux arrivent en grappes déclenchées par des événements communs, 400 par seconde, qui touchent les huit flux à la fois, comme une publication
macroéconomique ou une transaction importante sur un contrat à terme sur indice.

\begin{omfigure}
\begin{tikzpicture}
\begin{groupplot}[group style={group size=2 by 1,horizontal sep=1.7cm},width=6.6cm,height=5.2cm,
  tick label style={font=\small},label style={font=\small},grid=major,grid style={gray!25},table/col sep=comma]
\nextgroupplot[xmode=log,ymode=log,xmin=0.7,xmax=14000,ymin=3,ymax=300,xlabel={fenêtre (\unit{\micro\second}, log)},
  ylabel={débit de crête (\unit{\giga\bit\per\second}, log)},xtick={1,10,100,1000,10000},
  xticklabels={1,10,100,1\,000,10\,000},ytick={5,10,20,50,100,200},yticklabels={5,10,20,50,100,200}]
\addplot[omDef,very thick,mark=*,mark size=1.4pt] table[x=window_us,y=independent]{figdata/networks/01-networking-for-trading/peak.csv};
\addplot[omThm,very thick,mark=square*,mark size=1.4pt] table[x=window_us,y=common]{figdata/networks/01-networking-for-trading/peak.csv};
\addplot[gray,dashed,domain=0.7:14000] {10};
\nextgroupplot[xmode=log,xmin=8,xmax=6000,ymin=0,ymax=26,xlabel={mémoire tampon de sortie (Ko, log)},ylabel={trames perdues (\%)},
  xtick={10,100,1000},xticklabels={10,100,1\,000},
  legend columns=2,legend style={font=\footnotesize,at={(-0.2,-0.32)},anchor=north,draw=none,fill=none,
  /tikz/every even column/.append style={column sep=8pt}}]
\addplot[omDef,thick,dashed,mark=o,mark size=1.4pt] table[x=buffer_kb,y=f4_independent]{figdata/networks/01-networking-for-trading/fanin.csv};
\addplot[omThm,thick,dashed,mark=square,mark size=1.4pt] table[x=buffer_kb,y=f4_common]{figdata/networks/01-networking-for-trading/fanin.csv};
\addplot[omDef,very thick,mark=*,mark size=1.4pt] table[x=buffer_kb,y=f8_independent]{figdata/networks/01-networking-for-trading/fanin.csv};
\addplot[omThm,very thick,mark=square*,mark size=1.4pt] table[x=buffer_kb,y=f8_common]{figdata/networks/01-networking-for-trading/fanin.csv};
\legend{4 flux indépendants,4 flux avec événements communs,8 flux indépendants,8 flux avec événements communs}
\end{groupplot}
\end{tikzpicture}

\omcaption{Simulation~: flux vers un seul port de sortie de \qty{10}{\giga\bit\per\second}. À gauche~: la fenêtre la plus chargée de chaque longueur pour huit
flux (en pointillés~: le \omterm{def:nw:networking-for-trading:ser}{débit de ligne})~; les rafales indépendantes restent en dessous à partir de 100~\unit{\micro\second}, alors que les événements communs
le dépassent au-delà de \qty{60}{\giga\bit\per\second}. À droite~: trames perdues par rejet en queue selon la taille de la mémoire tampon, pour 50~ms de trafic.
La même charge moyenne nécessite environ 30~Ko de mémoire tampon lorsque les rafales sont indépendantes et environ un mégaoctet lorsqu'elles sont communes. Données~:
\texttt{fig\_bursts.py} sur \texttt{firm.netsim} (avec graine, déterministe).}\label{fig:nw:networking-for-trading:fanin}
\end{omfigure}

Les deux versions ont la même charge moyenne, et la différence tient tout entière à la corrélation
(\cref{fig:nw:networking-for-trading:fanin}). Avec des rafales indépendantes, les 100 microsecondes les plus chargées transportent
\qty{9.1}{\giga\bit\per\second}, et la mémoire tampon qui n'aurait rien perdu, soit le plus fort arriéré de la simulation, est de 21 à 31~Ko
sur vingt graines. Avec des événements communs, les 100 microsecondes les plus chargées transportent \qty{68}{\giga\bit\per\second}, et cette mémoire tampon
se situe entre 0.44 et 2.4~Mo sur les mêmes vingt graines, 1.1~Mo en moyenne. Retirer les événements communs constitue l'ablation~: le
mégaoctet est dû à la corrélation, et non à la charge. Les statistiques de trafic d'un port, moyennées sur une seconde, affichent la moitié d'un lien dans
les deux cas.

\begin{dated}{2026-09}{La mémoire tampon d'un commutateur à faible latence}\label{dat:nw:networking-for-trading:buffer}
La fiche technique de Cisco pour le Nexus 3548-X, un commutateur de 48 ports à \qty{10}{\giga\bit\per\second} vendu pour sa faible latence, indique sa
mémoire tampon comme «~6 Mo partagés entre 16 ports~; 18 Mo au total~» et des latences «~aussi basses que~» \qty{250}{\nano\second} en mode normal et
\qty{200}{\nano\second} en mode warp. Un port dont les quinze voisins sont également chargés peut compter sur environ un seizième de son bloc de
6~Mo, soit 375~Ko.
\end{dated}

Lu à côté de la simulation, l'encadré dit qu'un port qui peut emprunter tout le bloc partagé survit aux rafales d'événements communs
de la \cref{fig:nw:networking-for-trading:fanin}, et qu'un port qui n'obtient que son seizième n'y survit pas~: avec 500~Ko de mémoire tampon, les
huit flux avec événements communs perdent encore 2.6~\% de leurs trames. Les événements qui saturent tous les ports à la fois sont ceux qui
laissent à chaque port sa seule part. Les remèdes sont un port de sortie plus rapide (les rafales communes des huit flux atteignent
\qty{16}{\giga\bit\per\second} pendant une milliseconde, à l'intérieur d'un port de \qty{25}{\giga\bit\per\second}), moins de flux par port, ou un commutateur dont la
mémoire tampon est dédiée plutôt que partagée~; le chapitre~2 ajoute les équipements qui ne mettent rien en file d'attente.

\section{Tutoriel~: les microrafales, des trames aux pertes}\label{tut:nw:networking-for-trading:tut}

\begin{tutorial}
\textbf{Objectif.} Chiffrer les trames sur le fil, lire la courbe publiée des rafales et simuler un port de convergence jusqu'à la perte de trames.
\textbf{État final~:} le \cref{tab:nw:networking-for-trading:ser}, les \cref{fig:nw:networking-for-trading:opra,fig:nw:networking-for-trading:fanin},
et les mémoires tampons sans perte citées dans le texte.
\begin{tutsteps}
\item \textbf{Trames.} \texttt{nw\_wire.frame\_table()} et \texttt{mold\_example()} reproduisent le tableau et
  l'\cref{ex:nw:networking-for-trading:mold} à partir du \cref{lst:nw:networking-for-trading:frame}.
\item \textbf{La courbe publiée.} \texttt{opra\_curve()} convertit les pics de messages de l'opérateur en débits~; \texttt{opra\_capacity()}
  convertit l'avis de capacité en gigabits par seconde et en octets par paquet.
\item \textbf{Trafic.} \texttt{feed()} tire les paquets d'un flux à partir de \texttt{firm.netsim.bursty\_arrivals}, et
  \texttt{fanin\_traffic()} fusionne huit d'entre eux, avec ou sans événements communs (\cref{lst:nw:networking-for-trading:feed}).
\item \textbf{Le port.} \texttt{firm.netsim.egress} exécute la file à rejet en queue (\cref{lst:nw:networking-for-trading:egress})~;
  \texttt{fanin()} balaie la mémoire tampon, \texttt{zero\_loss\_buffer()} lance vingt graines avec une mémoire tampon non bornée et renvoie pour chaque
  exécution son plus grand arriéré. \texttt{python fig\_bursts.py} écrit les données des graphiques.
\end{tutsteps}
\textbf{Pour aller plus loin.} Faites varier la part des paquets accompagnés d'événements communs de 0 à 0.5 et tracez la mémoire tampon sans perte
en fonction de cette part~; remplacez le port de \qty{10}{\giga\bit\per\second} par un port de \qty{25}{\giga\bit\per\second} et trouvez la plus petite
mémoire tampon qui ne perd rien.
\end{tutorial}

\omcode{code/networks/01-networking-for-trading/python/nw_wire.py}{62}{77}{Un flux et une convergence de huit flux~: les instants des événements communs sont tirés une fois et
partagés par tous les flux.}\label{lst:nw:networking-for-trading:feed}

\omcode{code/firm/netsim/firm_netsim.py}{119}{140}{La file de sortie~: l'arriéré se résorbe au débit de ligne entre deux arrivées, et une
trame qui ne tient pas est rejetée entière.}\label{lst:nw:networking-for-trading:egress}

\section{Construction~: le modèle de réseau}\label{bld:nw:networking-for-trading:netsim}

\begin{build}
\bfield{Rôle} Un modèle au niveau du paquet du réseau de la société, que chaque chapitre ultérieur des parties~I et~II utilise~: trames, liaisons,
files, commutateurs et multicast, avec une comptabilité exacte des octets et un trafic amorcé.
\bfield{Interface} \texttt{firm\_netsim}~: constantes pour chaque en-tête~; \texttt{frame\_bytes}, \texttt{mold\_payload},
\texttt{wire\_bytes}, \texttt{ser\_ns}, \texttt{prop\_ns}~; \texttt{mcast\_mac}~; \texttt{peak\_rate(t, bytes, window)}~;
\texttt{event\_times}, \texttt{bursty\_arrivals}~; \texttt{egress(t, frames, gbps, buffer)} qui renvoie les départs, les rejets et
les arriérés~; \texttt{forward\_ns(frame, gbps, mode, fabric\_ns)} pour les équipements store-and-forward, cut-through et de couche~1~;
\texttt{replicate(groups, members, ports, snooping)}~; \texttt{merge}. Jumeau C++20 de la \omterm{def:nw:networking-for-trading:egress}{file de sortie},
\texttt{cpp/firm\_netsim.hpp}.
\bfield{Règles} Nanosecondes et octets entiers~; chaque trame paie 20 octets de préambule et d'intervalle~; les trames de moins de 64 octets sont
complétées~; le rejet en queue rejette les trames entières~; une file est FIFO~; chaque tirage aléatoire est amorcé~; la file C++ reproduit celle de Python à la
nanoseconde près.
\bfield{Tests d'acceptation} \texttt{code/firm/netsim/tests/}~: tailles de trames de paquets connus (64, 104, 1\,518 octets)~; correspondance d'adresses de la RFC~1112
et son recouvrement de 32 vers 1~; la fenêtre de débit de crête sur un flux construit à la main~; ordre FIFO, conservation des trames et
la borne de la mémoire tampon de la \omterm{def:nw:networking-for-trading:egress}{file de sortie}~; dix trames arrivant ensemble quittent après dix sérialisations~; l'échantillon partagé avec le
test C++.
\bfield{Pour aller plus loin} Files à priorité (priorité stricte aux ordres sur les données de marché sur une liaison montante partagée)~; un commutateur à mémoire partagée dans lequel
les ports empruntent à un réservoir commun selon un seuil dynamique, pour reproduire l'arithmétique de l'\cref{dat:nw:networking-for-trading:buffer}
en simulation.
\end{build}

\begin{omsources}
\item OPRA, \emph{Key Operating Metrics of U.S. Options Securities Information Processor} (août 2026)~; SIAC, \emph{Revised
OPRA Capacity Projections} (15 septembre 2025).
\item G.~Fairhurst, «~Ethernet Frame Calculations~» (University of Aberdeen)~; UNH InterOperability Laboratory, \emph{10
Gigabit Ethernet} tutoriel MAC.
\item RFC~768 (UDP), RFC~791 (IPv4), RFC~9293 (TCP), RFC~896 (Nagle), RFC~1122 (exigences pour les hôtes), RFC~1112 et RFC~5771
(adressage multicast), RFC~3376 (IGMPv3), RFC~4541 (snooping)~; Linux \texttt{tcp(7)}.
\item Cisco, \emph{Nexus 3548-X, 3524-X, 3548-XL and 3524-XL Switches Data Sheet}.
\end{omsources}

\section{Exercices}

\begin{exercise}[$\star$]\label{exo:nw:networking-for-trading:1}
Calculer le \omterm{def:nw:networking-for-trading:ser}{délai de sérialisation} d'une trame de 256 octets à \qty{25}{\giga\bit\per\second}, et le nombre de trames de 64 octets qu'une
liaison de \qty{10}{\giga\bit\per\second} peut transporter en une seconde~?
\end{exercise}

\begin{exercise}[$\star$]\label{exo:nw:networking-for-trading:2}
À quelle adresse Ethernet le \omterm{def:nw:networking-for-trading:group}{groupe multicast} 233.54.12.111 est-il associé~? Donner un autre groupe de 224.0.0.0/4 qui partage cette adresse.
\end{exercise}

\begin{exercise}[$\star$]\label{exo:nw:networking-for-trading:3}
D'après l'\cref{dat:nw:networking-for-trading:opra}, quel débit une société qui reçoit les deux copies du flux, avec la
marge de retransmission, doit-elle prévoir pour la crête de 10 millisecondes, et combien de liaisons de \qty{25}{\giga\bit\per\second} cela représente-t-il~?
\end{exercise}

\begin{exercise}[$\star\star$]\label{exo:nw:networking-for-trading:4}
Une passerelle écrit un nouvel ordre puis, quarante microsecondes plus tard, une annulation d'un autre ordre sur la même connexion TCP, sans
\texttt{TCP\_NODELAY}. La plateforme n'acquitte les données qu'avec ses propres messages ou à l'échéance de son minuteur d'acquittement différé. Que
peut-il arriver à l'annulation, et pourquoi l'ordre des deux n'est-il jamais menacé~?
\end{exercise}

\begin{exercise}[$\star\star$]\label{exo:nw:networking-for-trading:5}
Quatre entrées émettent chacune à \qty{6}{\giga\bit\per\second} pendant 50 microsecondes vers un port de \qty{10}{\giga\bit\per\second} dont la
file est vide. Combien d'octets sont en file à la fin de la rafale, et combien de temps attend le dernier octet~?
\end{exercise}

\begin{exercise}[$\star\star$]\label{exo:nw:networking-for-trading:6}
Dans la \cref{fig:nw:networking-for-trading:opra}, calculer, en janvier 2024 et en juillet 2026, le rapport entre le débit de la milliseconde la plus
chargée et celui de la seconde la plus chargée. Qu'est-ce qui a changé, et qu'est-ce que cela implique pour une liaison dimensionnée en 2024~?
\end{exercise}

\begin{exercise}[$\star\star\star$]\label{exo:nw:networking-for-trading:7}
\emph{Programmation.} Avec \texttt{nw\_wire.zero\_loss\_buffer}, calculer la mémoire tampon sans perte moyenne sur vingt graines lorsque 10~\% des paquets
arrivent avec des événements communs, puis lorsque 50~\% des paquets y arrivent. Comment la mémoire tampon évolue-t-elle avec cette part~?
\end{exercise}

\begin{exercise}[$\star\star\star$]\label{exo:nw:networking-for-trading:8}
\emph{Trouver la faille.} «~Notre port de données de marché moyenne \qty{3}{\giga\bit\per\second} sur une liaison de \qty{10}{\giga\bit\per\second} et
son pic sur une minute est de \qty{4.2}{\giga\bit\per\second}. Plus de la moitié de la liaison nous reste disponible~; les pertes de paquets que signale notre handler
doivent venir du serveur.~»
\end{exercise}

\section{Problème~: L'ouverture qui ne tenait pas}

\begin{problem}[{Problème du week-end --- huit flux, un port et un tampon partagé}]\label{pb:nw:networking-for-trading:1}
Le commutateur de données de marché d'une société envoie huit flux de plateformes vers un serveur de stratégie par un même port de \qty{10}{\giga\bit\per\second}. Chaque
flux a un débit moyen de \qty{0.6}{\giga\bit\per\second}~; lors d'un événement touchant l'ensemble du marché, les huit flux émettent ensemble en rafale à \qty{2.5}{\giga\bit\per\second}
chacun pendant une milliseconde. Les trames font en moyenne 200 octets. Le commutateur est celui de l'\cref{dat:nw:networking-for-trading:buffer}~; pendant
un tel événement, chaque port du bloc est occupé, de sorte que le port dispose du seizième du tampon partagé de 6~Mo.

\medskip\noindent\textbf{Partie I --- Moyennes.}
\begin{enumerate}
\item Quelle est la charge moyenne du port et son taux d'utilisation~?
\item Combien de trames par seconde transporte-t-il en moyenne (préambule et intervalle inter-trames compris)~?
\item Quel est le débit offert pendant l'événement, et le taux d'utilisation correspondant~?
\item Pourquoi le compteur d'une minute du port ne fait-il jamais apparaître l'événement~?
\end{enumerate}

\medskip\noindent\textbf{Partie II --- La file d'attente.}
\begin{enumerate}[resume]
\item Combien d'octets sont en attente en file à la fin de la rafale, avec un tampon illimité~?
\item Combien de trames cela représente-t-il~?
\item Combien de temps le dernier octet de la rafale attend-il dans la file~?
\item Après la rafale, les flux reviennent à leurs moyennes. Combien de temps la file met-elle à se vider~?
\end{enumerate}

\medskip\noindent\textbf{Partie III --- Le tampon.}
\begin{enumerate}[resume]
\item Quel tampon ne perdrait rien~?
\item Quel tampon le port obtient-il réellement pendant l'événement~?
\item Combien d'octets et combien de trames sont écartés~?
\item Quelle part des trames de la rafale cela représente-t-elle~?
\item Le bloc entier de 6~Mo aurait-il suffi~?
\end{enumerate}

\medskip\noindent\textbf{Partie IV --- Le verdict.}
\begin{enumerate}[resume]
\item Énoncer le \emph{résultat nommé}~: le tampon minimal pour une perte nulle, et le nombre de trames écartées avec la part du port.
\item Combien de temps la rafale devrait-elle durer pour que le bloc entier de 6~Mo déborde~?
\item Un port de \qty{25}{\giga\bit\per\second} mettrait-il seulement des trames en file pendant l'événement~?
\item La société prend les deux copies de chaque flux sur des ports distincts du même bloc. Cela aide-t-il~?
\item Quel compteur du commutateur montre la perte, et quel compteur du gestionnaire~?
\item Comment les chemins de reprise du gestionnaire (livre~13, chapitre~18) gèrent-ils la perte~?
\item En une phrase~: pourquoi l'événement qui remplit un port est-il celui qui remplit tous les ports~?
\end{enumerate}
\end{problem}

\section{Questions d'entretien}

\begin{interviewq}[$\star$ \iqroles{développeur}]\label{iq:nw:networking-for-trading:1}
Qu'est-ce que le \omterm{def:nw:networking-for-trading:ser}{délai de sérialisation}~? Pourquoi passer de 10 à \qty{25}{\giga\bit\per\second} compte-t-il pour un ordre qui tient dans une petite
trame~?
\end{interviewq}

\begin{interviewq}[$\star$ \iqroles{développeur, trader}]\label{iq:nw:networking-for-trading:2}
Pourquoi les données de marché sont-elles envoyées en multicast UDP, et les ordres en TCP~?
\end{interviewq}

\begin{interviewq}[$\star\star$ \iqroles{développeur}]\label{iq:nw:networking-for-trading:3}
Expliquez l'\omterm{def:nw:networking-for-trading:nagle}{algorithme de Nagle} et l'acquittement différé. Que configurez-vous sur une socket de saisie des ordres, et pourquoi~?
\end{interviewq}

\begin{interviewq}[$\star\star$ \iqroles{développeur}]\label{iq:nw:networking-for-trading:4}
Le feed handler signale des trous sur les deux lignes à l'ouverture, mais la carte du serveur n'affiche aucune perte. Où regardez-vous~?
\end{interviewq}

\begin{interviewq}[$\star\star$ \iqroles{développeur}]\label{iq:nw:networking-for-trading:5}
Que fait \omterm{def:nw:networking-for-trading:snooping}{IGMP snooping}, et que se passe-t-il dans un réseau de trading qui en est dépourvu~?
\end{interviewq}

\begin{interviewq}[$\star\star\star$ \iqroles{développeur, chercheur}]\label{iq:nw:networking-for-trading:6}
Comment dimensionneriez-vous le tampon d'un port de commutateur qui reçoit plusieurs flux de données de marché~?
\end{interviewq}
